\chapter{Opis implementacji}

W niniejszym rozdziale przedstawiono sposób implementacji rozwiązania opisanego od strony teoretycznej w poprzednim rozdziale. Implementacja obejmuje generowanie grafów, trzy warianty algorytmu Dijkstry, analizę otrzymanych wyników oraz towrzenie wykresów umożliwiających ich porównanie.

Poszczególne etapy działania programu zostały rozdzielone, żeby mogły być uruchamiane niezależnie. Dane pomiędzy nimi przekazywane są przy użyciu plików \texttt{JSON}. Wygenerowane grafy są zapisywane do plików, a następnie odczytywane przez moduły wykonujące odpowiednie wersje algorytmu Dijkstry. Wyniki zapisywane są w plikach \texttt{JSON}, które następnie wykorzystywane są przez moduły do analizy statystycznej. Jej wyniki stanowią z kolei dane wejściowe dla modułów odpowiedzialnych za generowanie wykresów. Cały proces przetwarzania danych w implementacji można przedstawić w następujący sposób:
\begin{center}
    generowanie grafów
    $\rightarrow$
    zapis do \texttt{JSON}
    $\rightarrow$
    wykonanie algorytmów
    $\rightarrow$
    zapis wyników do \texttt{JSON}
    $\rightarrow$
    analiza statystyczna
    $\rightarrow$
    zapis statystyk do \texttt{JSON}
    $\rightarrow$
    generowanie wykresów.
\end{center}

Taki podział na niezależne moduły ułatwia przeprowadzanie eksperymentów z wykorzystaniem poszczególnych wersji algorytmu Dijkstry, analizę uzyskanych wyników oraz ponowne generowanie wykresów bez konieczności wykonywania całego procesu od początku.
\newpage

\section{Tworzenie grafów}

Pierwszym etapem jest przygotowanie grafów stanowiących dane wejściowe dla wszystkich analizowanych wariantów algorytmu. W części teoretycznej przedstawiono cztery rodzaje grafów wykorzystywane w pracy: graf rzadki, graf o średniej gęstości, graf gęsty oraz przypadek szczególny. Ich właściwości, a także przyjęte liczby krawędzi, opisano w rozdziale dotyczącym rodzajów grafów~\ref{sec:graph_types}.

Liczba wierzchołków grafów oraz liczba różnych grafów generowanych dla każdego badanego rozmiaru są określane za pomocą pliku konfiguracyjnego. Pozwala to na łatwe dostosowanie parametrów bez konieczności modyfikowania kodu źródłowego. 

Grafy przechowywane są w postaci listy sąsiedztwa, zgodnie z reprezentacją opisaną wcześniej w rozdziale~\ref{sec:graph_representation}. Po wygenerowaniu każdy graf jest zapisywany do pliku  \texttt{JSON}. Jego nazwa zawiera informacje pozwalające określić typ grafu, jego rozmiar oraz numer. Dzięki temu wszystkie implementacje algorytmu Dijkstry mogą korzystać z tych samych danych wejściowych.


\subsection{Generowanie grafu gęstego}

Graf gęsty został zaimplementowany jako graf pełny. Oznacza to, że dla każdej pary różnych wierzchołków tworzona jest dokładnie jedna krawędź. Liczba krawędzi jest zatem zgodna z zależnością

\begin{equation}
    |E| = \frac{|V|(|V|-1)}{2},
\end{equation}
\noindent
przedstawioną również w części teoretycznej w rozdziale~\ref{subsec:graph_dense}. Dla każdej utworzonej krawędzi losowana jest jej waga. Przykład grafu gęstego przedstawiono na rys.~\ref{fig:graph_dense}.
\newpage

\subsection{Generowanie grafu o średniej gęstości}

W przypadku grafu o średniej gęstości liczba tworzonych krawędzi stanowi połowę maksymalnej liczby krawędzi grafu pełnego, zgodnie z zależnością przedstawioną wcześniej w rozdziale~\ref{subsec:graph_half_dense}

\begin{equation}
    |E| = \frac{|V|(|V|-1)}{4}.
\end{equation}

Kolejne krawędzie tworzone są pomiędzy losowo wybieranymi parami różnych wierzchołków. Podczas generowania sprawdzane jest, czy dana krawędź nie została utworzona wcześniej, dzięki czemu graf nie zawiera wielokrotnych krawędzi pomiędzy tą samą parą wierzchołków. Wagi krawędzi losowane są z zakresu określonego w części teoretycznej. Przykład wykorzystywanego grafu o średniej gęstości przedstawiono na rys.~\ref{fig:graph_half_edges}.

\subsection{Generowanie grafu rzadkiego}

Graf rzadki generowany jest zgodnie z definicją przyjętą w części teoretycznej w rozdziale~\ref{subsec:graph_sparse}, zgodnie z którą liczba krawędzi jest równa liczbie wierzchołków

\begin{equation}
    |E| = |V|.
\end{equation}
\noindent
Mechanizm losowania krawędzi oraz ich wag jest analogiczny jak w przypadku grafu o średniej gęstości. Różnica polega na większej liczbie tworzonych połączeń. Przykład grafu tego typu przedstawiono na rys.~\ref{fig:graph_sparse}.


\subsection{Generowanie przypadku szczególnego}

Ostatnim generowanym typem jest przypadek szczególny przedstawiony na rys.~\ref{fig:graph_special_case}. Podobnie jak graf gęsty jest on grafem pełnym, jednak wagi jego krawędzi nie są losowane. Sposób ich wyznaczania odpowiada zależności przedstawionej w części teoretycznej w rozdziale~\ref{subsec:graph_special_case}. 
\[
w(i,j)=
\begin{cases}
1, & \text{gdy } j=i+1,\\[6pt]
2\left(|V|-i\right), & \text{gdy } j>i+1.
\end{cases}
\]

\section{Implementacja algorytmu Dijkstry}
\label{sec:dijkstra_implementation}

Podstawowy sposób działania algorytmu Dijkstry został przedstawiony w alg.~\ref{alg:dijkstra}. Wszystkie trzy implementacje użyte w części działają według tego samego schematu. Różnią się przede wszystkim sposobem wykonania operacji wyboru nieodwiedzonego wierzchołka o najmniejszej aktualnie znanej odległości.

W celu porównania różnych metod wyszukiwania minimum zaimplementowano trzy warianty:
\begin{itemize}
    \item wersję naiwną,
    \item wersję wykorzystującą kopiec binarny,
    \item wersję wykorzystującą kwantowy algorytm wyszukiwania minimum.
\end{itemize}

\subsection{Wersja naiwna}
\label{subsec:naive_implementation}

W wersji naiwnej wyszukiwanie minimum realizowane poprzez liniowe przeglądanie wszystkich wierzchołków grafu. W każdej iteracji algorytmu Dijkstry wybierany jest spośród nieprzetworzonych wierzchołków ten, którego aktualnie znana odległość od wierzchołka źródłowego jest najmniejsza. Sposób ten odpowiada wersji naiwnej opisanej w części teoretycznej w rozdziale~\ref{subsec:naive_minimum}.

Na potrzeby zliczania operacji podczas wyszukiwania wprowadzono zmienną przechowującą liczbę wykonanych operacji. Za pojedynczą operację przyjęto sprawdzenie jednego wierzchołka w wewnętrznej pętli wyszukiwania. Licznik ten jest zwiększany o $1$ dla każdego sprawdzanego wierzchołka, niezależnie od tego, czy zostanie on nowym kandydatem na minimum. Ponieważ podczas pojedynczego wyszukiwania przeglądane jest zawsze $n$ wierzchołków, liczba zliczonych operacji wynosi

\begin{equation}
    C_{\mathrm{search}} = n.
\end{equation}

Wyszukiwanie wykonywane jest $n$ razy, ponieważ zewnętrzna pętla algorytmu również wykonuje $n$ iteracji. Całkowita liczba operacji zliczonych dla jednego uruchomienia wersji naiwnej wynosi więc

\begin{equation}
    C_{\mathrm{naive}}
    =
    n^2.
\end{equation}

Po zakończeniu działania algorytmu wartość zmiennej liczącej liczbe operacji jest zapisywana, a następnie przekazywana w pliku \texttt{JSON} do dalszej analizy statystycznej.


\subsection{Wersja z kopcem binarnym}
\label{subsec:binary_implementation}

Druga klasyczna implementacja wykorzystuje kopiec binarny opisany w części teoretycznej w rozdziale~\ref{subsec:naive_minimum}. Przykładową strukturę kopca przedstawiono na rys.~\ref{fig:binary_heap_list}. W przeciwieństwie do wersji naiwnej wybór wierzchołka o najmniejszej odległości nie wymaga liniowego przeglądania wszystkich nieodwiedzonych wierzchołków.

Na potrzeby algorytmu zaimplementowano własną strukturę kopca, która zlicza wykonywane podczas jego działania operacje. Uwzględniane są ich dwa rodzaje:
\begin{itemize}
\item porównania wartości priorytetów,
\item zamiany elementów wewnątrz kopca.
\end{itemize}
\noindent
Operacje te są zliczane podczas wstawiania elementów do kopca oraz pobierania elementu o najmniejszym priorytecie. 

W implementacji zastosowano podejście \textit{lazy heap}. Początkowo do kopca wstawiany jest wyłącznie wierzchołek startowy z odległością równą zero. Pozostałe wierzchołki są dodawane dopiero w momencie znalezienia do nich pierwszej skończonej odległości. W przypadku znalezienia krótszej drogi podczas kolejnej relaksacji wartość przechowywana wcześniej w kopcu nie jest bezpośrednio modyfikowana za pomocą operacji \textit{decrease-key}. Zamiast tego aktualizowana jest odległość danego wierzchołka, a do kopca dodawany jest nowy element zawierający tę odległość.

Wstawienie nowego elementu może wymagać ponownego uporządkowania kopca, co wiąże się z kolejnymi porównaniami i zamianami elementów. Takie same operacje mogą wystąpić podczas pobierania elementu o najmniejszym priorytecie. Dlatego liczba operacji dla wersji z kopcem jest sumą liczników obu operacji. Po zakończeniu algorytmu suma ta jest zapisywana do pliku \texttt{JSON}.
\newpage

\subsection{Wersja kwantowa}
\label{subsec:quantum_implementation}

Trzecia implementacja algorytmu Dijkstry wykorzystuje kwantowy algorytm wyszukiwania minimum Dürra-Høyera przedstawiony w alg.~\ref{alg:quantum_minimum}. Algorytm ten wykorzystuje algorytm BBHT (alg.~\ref{alg:bbht}), który z kolei wykonuje określoną liczbę iteracji algorytmu Grovera (alg.~\ref{alg:grover}).

Ogólna struktura algorytmu Dijkstry pozostaje taka sama jak w alg.~\ref{alg:dijkstra}. W każdej iteracji tworzona jest lista
wierzchołków wraz z odpowiadającymi im odległościami. Lista może zawierać również wierzchołki o odległości równej nieskończoności. Takie wierzchołki mogą zostać pominięte podczas dostosowywania rozmiaru przestrzeni wyszukiwania dlatego liczba elementów jest następnie zwiększana lub zmniejszana do najbliższej odpowiedniej potęgi liczby~2. Wynika to ze sposobu reprezentacji stanów przez rejestr kwantowy. Rejestr składający się z $n$ kubitów pozwala reprezentować $2^n$ stanów bazowych. Liczba operacji w wersji kwantowej wyznaczana jest na podstawie dwóch elementów:
\begin{itemize}
    \item kosztu przygotowania rejestru kwantowego w stanie równomiernej superpozycji,
    \item liczby wykonanych iteracji algorytmu Grovera.
\end{itemize}
\noindent
Dla przestrzeni wyszukiwania zawierającej $N$ elementów potrzebny jest rejestr składający się z
\begin{equation}
    n = \log_2 N
\end{equation}
\noindent
kubitów, ponieważ $N=2^n$. Przygotowanie superpozycji stanów wymaga zastosowania bramki Hadamarda do każdego z $n$ kubitów, zgodnie ze
wzorem~\eqref{eq:superposition}. Koszt przygotowania rejestru wynosi zatem
\begin{equation}
    C_{\mathrm{init}} = \log_2 N.
\end{equation}

W procedurze BBHT losowana jest dodatkowo wartość $j$, określająca liczbę iteracji algorytmu Grovera wykonywanych w danej próbie. Koszt pojedynczej próby BBHT określono więc jako
\begin{equation}
    C_{\mathrm{BBHT}}
    =
    \log_2 N + j.
\end{equation}

Pierwszy składnik odpowiada koszcie przygotowania superpozycji, natomiast drugi liczbie wykonanych iteracji algorytmu Grovera. Jeżeli algorytm BBHT wymaga kilku prób,
ich koszty są sumowane. Algorytm Dürra-Høyera może następnie wielokrotnie wywoływać algorytm BBHT podczas znajdowania kolejnych elementów mniejszych od aktualnego
kandydata na minimum. Koszty tych wywołań są również sumowane. Ostatecznie liczba operacji zarejestrowana podczas działania algorytmu Dijkstry odpowiada sumie kosztów wszystkich prób BBHT wykonanych podczas wszystkich wywołań kwantowego wyszukiwania minimum.

Dodatkowo rejestrowane są: liczba wywołań algorytmu wyszukiwania minimum, liczba przypadków, w których minimum znalezione przez algorytm kwantowy różniło się od rzeczywistego minimum, oraz informacja o niepoprawności końcowego wyniku algorytmu Dijkstry. Po zakończeniu działania wszystkie zliczone wartości zapisywane są do pliku w formacie \texttt{JSON}.

\section{Analiza wyników}
\label{sec:results_analysis_implementation}

Na tym etapie działania programu dane są wczytywane z plików \texttt{JSON} zapisanymi w poprzednim etapie opisanym w \ref{sec:dijkstra_implementation}, a następnie poddawane analizie. Takie rozwiązanie pozwala na modyfikowanie sposobu obliczania statystyk bez konieczności ponownego uruchamiania wcześniejszych etapów programu.

\subsection{Ogólna analiza wyników}
\label{subsec:broad_results_analysis}

Dla każdej grupy pomiarów, obejmującej naiwną implementację klasycznego algorytmu Dijkstry, implementację wykorzystującą kopiec binarny oraz kwantowy algorytm Dijkstry, obliczane są następujące statystyki dotyczące liczby operacji wykonywanych przez poszczególne algorytmy:
\begin{itemize}
\item średnia,
\item mediana,
\item odchylenie standardowe,
\item wartość minimalna i maksymalna,
\end{itemize}
\noindent
W dalszej analizie wyników szczególnie istotne są mediana oraz odchylenie standardowe. Mediana jest wykorzystywana, ponieważ jest mniej wrażliwa na pojedyncze wartości odstające niż średnia. Odchylenie standardowe pozwala natomiast określić sam rozrzut wyników. Ma to znaczenie przede wszystkim w przypadku wersji kwantowej, ale jest obliczane też dla obu implementacji klasycznych, dzięki czemu możliwe jest porównanie stabilności wszystkich analizowanych wersji algorytmu.

Obliczone statystyki zapisywane są do kolejnych plików \texttt{JSON}, które stanowią dane wejściowe dla modułów odpowiedzialnych za tworzenie wykresów.

\subsection{Analiza wyników wersji kwantowej}
\label{subsec:quantum_results_analysis}

W przypadku wersji kwantowej przeprowadzana jest dodatkowa analiza związana z probabilistycznym działaniem algorytmu wyszukiwania minimum Dürra-Høyera przedstawionego w algorytmie~\ref{alg:quantum_minimum}. Algorytm ten wykorzystuje algorytm BBHT (alg.~\ref{alg:bbht}), która wykonuje określoną liczbę iteracji algorytmu Grovera (alg.~\ref{alg:grover}).

Oprócz analizy liczby wykonanych operacji opisanej w  poprzednim podrozdziale~\ref{subsec:broad_results_analysis}, analizowane są również: liczba wywołań algorytmu wyszukiwania minimum, liczba przypadków, w których znalezione minimum różniło się od opprawnego minimum, oraz liczba niepoprawnych wyników samego algorytmu Dijkstry. Dla tych wartości obliczane są podstawowe statystyki, takie jak średnia, mediana, odchylenie standardowe, minimum oraz maksimum. Porównywana jest również liczba błędnie znalezionych minimów z liczbą niepoprawnych wyników algorytmu Dijkstry. Pozwala to sprawdzić, czy błędne znalezienie minimum zawsze prowadzi do niepoprawnego wyniku  algorytmu Dijkstry.

Dodatkowo analizowana jest skuteczność algorytmu Dijkstry, określając liczbę poprawnych wyników wśród wszystkich otrzymanych wyników. W podobny sposób oceniana jest skuteczność algorytmu wyszukiwania minimum Dürra-Høyera na podstawie tego, jak często znalezione minimum było zgodne z poprawnym minimum.

Oddzielnie analizowane są konfiguracje z limitem wynikającym ze wzoru~\eqref{eq:minimum_time} oraz bez limitu. Pozwala to sprawdzić, jaki wpływ zastosowanie limitu ma na liczbę wykonywanych operacji oraz na prawdopodobieństwo otrzymania poprawnego wyniku.

Wszystkie wyniki przeprowadzonej analizy są następnie zapisywane do plików \texttt{JSON}.
\newpage

\subsection{Tworzenie wykresów}
\label{subsec:plots_implementation}

Ostatnim etapem programu jest przygotowanie wykresów przedstawiających uzyskane wyniki. Dane potrzebne do ich utworzenia są odczytywane z plików  \texttt{JSON}, które powstają podczas wcześniejszej analizy wyników opisanej w rozdziale~\ref{sec:results_analysis_implementation}. Rozwiązanie to pozwala zmieniać sposób prezentacji danych bez konieczności ponownego wykonywania wcześniejszych etapów programu.

Dla wszystkich implementacji tworzone są wykresy przedstawiające średnią, medianę oraz odchylenie standardowe w zależności od liczby wierzchołków grafu. Pozwala to obserwować, jak wraz ze wzrostem rozmiaru grafu zmienia się liczba operacji.

Dla wersji kwantowej tworzone są również dodatkowe wykresy. Przedstawiają one liczbę błędnie znalezionych minimów, liczbę niepoprawnych wyników algorytmu Dijkstry, liczbę wywołań algorytmu wyszukiwania minimum Dürra-Høyera oraz skuteczność jego działania.






